\documentclass[11pt,a4paper]{article}

\usepackage[a4paper,margin=1.8cm]{geometry}
\usepackage{fontspec}
\usepackage{graphicx}
\usepackage[export]{adjustbox}
\usepackage{xcolor}
\usepackage{titlesec}
\usepackage{enumitem}
\usepackage{fancyhdr}
\usepackage{hyperref}
\usepackage{caption}
\usepackage{float}
\usepackage{tabularx}
\usepackage{array}
\usepackage{booktabs}
\usepackage{needspace}
\usepackage{amssymb}

\setmainfont{Arial}
\definecolor{icsred}{HTML}{C91624}
\definecolor{icsblack}{HTML}{111111}
\definecolor{icsgray}{HTML}{5F5A55}
\definecolor{icslight}{HTML}{F4F1EA}
\definecolor{icsnote}{HTML}{E9E9E9}

\hypersetup{
  colorlinks=true,
  linkcolor=icsred,
  urlcolor=icsred
}

\pagestyle{fancy}
\setlength{\headheight}{14pt}
\fancyhf{}
\lhead{\textcolor{icsgray}{Panduan Admin HR Assistant}}
\rhead{\textcolor{icsgray}{ICS Compute}}
\cfoot{\textcolor{icsgray}{\thepage}}
\renewcommand{\headrulewidth}{0.3pt}

\titleformat{\section}
  {\Large\bfseries\color{icsblack}}
  {\thesection}{0.75em}{}
\titleformat{\subsection}
  {\normalsize\bfseries\color{icsblack}}
  {\thesubsection}{0.75em}{}
\titlespacing*{\section}{0pt}{1.1em}{0.55em}
\titlespacing*{\subsection}{0pt}{0.9em}{0.35em}

\setlist[itemize]{leftmargin=1.2em,itemsep=0.2em,topsep=0.25em}
\setlist[enumerate]{leftmargin=1.4em,itemsep=0.2em,topsep=0.25em}
\captionsetup{font=small,labelfont=bf}
\setlength{\parindent}{0pt}
\setlength{\parskip}{0.3em}
\newcolumntype{Y}{>{\raggedright\arraybackslash}X}
\newcolumntype{B}[1]{>{\bfseries\raggedright\arraybackslash}p{#1}}
\newcommand{\checkbox}{\textcolor{icsred}{\checkmark}}
\newcommand{\checkline}[1]{\checkbox & #1 \\[0.5em]}

\newcommand{\screenshot}[3][0.72\linewidth]{
  \begin{figure}[H]
    \centering
    \includegraphics[max width=#1,max height=0.46\textheight,keepaspectratio]{assets/#2}
    \caption{#3}
  \end{figure}
}

\newcommand{\notebox}[1]{
  \vspace{0.25em}
  \noindent
  \colorbox{icsnote}{
    \parbox{\dimexpr\linewidth-2\fboxsep-10pt}{
      \small
      \textbf{Catatan:} #1
    }
  }
  \vspace{0.35em}
}

\begin{document}

\begin{titlepage}
  \centering
  \vspace*{2.2cm}
  {\Huge\bfseries How To Use Admin\par}
  \vspace{0.4cm}
  {\Huge\bfseries \textcolor{icsred}{HR Assistant}\par}
  \vspace{0.9cm}
  {\Large ICS Compute\par}
  \vfill
  \begin{tabular}{>{\bfseries}rl}
    Versi & 2.0 \\
    Tanggal & 3 Agustus 2026 \\
    Akses & Internal HR / Admin \\
  \end{tabular}
  \vfill
  {\large Dokumen panduan operasional admin\par}
\end{titlepage}

\tableofcontents
\newpage

\Needspace{10\baselineskip}
\section{Akses Aplikasi}

Buka HR Assistant melalui URL berikut:

\begin{center}
\texttt{http://ec2-98-81-149-80.compute-1.amazonaws.com/}
\end{center}

\begin{itemize}
  \item Akses aplikasi harus menggunakan WiFi kantor.
  \item Gunakan awalan \texttt{http://}, bukan \texttt{https://}.
  \item Login memakai akun Google Workspace \texttt{@icscompute.com}. Akun Gmail
        pribadi atau domain lain akan ditolak oleh sistem.
\end{itemize}

\Needspace{20\baselineskip}
\section{Login dengan Google}

Aplikasi tertutup untuk publik: sebelum login, seluruh layar tertutup halaman
sign-in dan tidak ada menu yang bisa dibuka. Login tidak lagi memakai email dan
password; semuanya lewat akun Google Workspace kantor.

\screenshot[0.66\linewidth]{00-login-entry.png}{Halaman sign-in yang menutup aplikasi sebelum login.}

\begin{enumerate}
  \item Klik tombol \textbf{Sign in with Google}.
  \item Pilih akun Google Workspace \texttt{@icscompute.com} Anda.
  \item Aplikasi terbuka otomatis setelah akun terverifikasi.
\end{enumerate}

\notebox{Jika browser Anda sudah punya sesi Google aktif, akan muncul pop-up
\textit{One Tap} di kanan atas dan Anda bisa langsung memilih akun tanpa menekan
tombol. Jika pop-up tidak muncul (mis. diblokir browser), gunakan tombol
\textbf{Sign in with Google} pada kartu tersebut.}

\Needspace{12\baselineskip}
\subsection{Perbedaan User dan Admin}

Semua karyawan dengan email \texttt{@icscompute.com} bisa login sebagai
\textbf{user} dan memakai chat, FAQ, serta Docs. Menu pengelolaan hanya muncul
jika email Anda terdaftar sebagai \textbf{admin}.

\begin{center}
\renewcommand{\arraystretch}{1.2}
\begin{tabularx}{0.94\linewidth}{@{}B{3.4cm}Y@{}}
\toprule
Role & Yang bisa diakses \\
\midrule
User & Chat, FAQ, Docs (unduh dokumen dan form). \\
Admin & Semua akses user, ditambah tombol \textbf{New admin}, menu
\textbf{Logs}, \textbf{Analytics}, \textbf{Guardrails}, serta panel admin di FAQ
dan Docs. \\
\bottomrule
\end{tabularx}
\end{center}

Setelah login sebagai admin, sidebar menampilkan enam menu dan nama Anda muncul
di panel akun kiri bawah dengan label \textbf{Admin}.

\screenshot[0.76\linewidth]{02-admin-mode.png}{Tampilan setelah login sebagai admin.}

\notebox{Admin pertama ditentukan lewat konfigurasi server
(\texttt{INITIAL\_ADMIN\_EMAIL}). Admin berikutnya ditambahkan dari dalam
aplikasi seperti pada bab berikut.}

\Needspace{16\baselineskip}
\section{Pilihan Bahasa (EN / ID)}

Seluruh tampilan tersedia dalam Bahasa Inggris dan Bahasa Indonesia. Tombol
\textbf{EN / ID} ada di kanan atas, dan juga di kartu sign-in sehingga bisa
diganti sebelum login.

\screenshot[0.86\linewidth]{21-language-toggle.png}{Tombol pilihan bahasa pada bar atas.}

\begin{itemize}
  \item Klik \textbf{ID} untuk Bahasa Indonesia, \textbf{EN} untuk Bahasa Inggris.
  \item Pilihan bahasa tersimpan di browser, jadi tetap sama saat aplikasi dibuka lagi.
  \item Pilihan ini hanya mengubah tampilan aplikasi, bukan bahasa jawaban chatbot.
        Chatbot menjawab mengikuti bahasa pertanyaan Anda.
\end{itemize}

\screenshot[0.76\linewidth]{22-language-indonesian.png}{Tampilan aplikasi dalam Bahasa Indonesia.}

\Needspace{16\baselineskip}
\section{Menambahkan Admin Baru}

Admin dapat memberi akses admin ke rekan lain lewat tombol \textbf{New admin}.
Formulir ini hanya meminta email, karena password sepenuhnya ditangani Google.

\begin{enumerate}
  \item Klik \textbf{New admin} di kanan atas.
  \item Isi email Google Workspace orang tersebut, harus berakhiran
        \texttt{@icscompute.com}.
  \item Klik \textbf{Save admin}.
\end{enumerate}

\screenshot[0.46\linewidth]{03-new-admin.png}{Form tambah admin baru.}

\notebox{Akses admin berlaku untuk email persis seperti yang diisi, dan langsung
aktif saat orang tersebut login dengan akun Google-nya. Jika email sudah
terdaftar sebagai admin, sistem akan menolak penambahan.}

\Needspace{16\baselineskip}
\section{Chat dan Tombol Stop}

Admin tetap dapat menggunakan chat seperti user biasa. Saat chatbot sedang
memproses jawaban, tombol kirim berubah menjadi tombol \textbf{Stop}.

\begin{enumerate}
  \item Tulis pertanyaan pada kolom chat.
  \item Klik tombol kirim.
  \item Jika ingin menghentikan proses, klik tombol \textbf{Stop}.
\end{enumerate}

\screenshot[0.76\linewidth]{04-chat-stop.png}{Chat sedang memproses jawaban dan dapat dihentikan dengan tombol Stop.}

\notebox{Chatbot menjawab berdasarkan SOP/knowledge document yang sudah masuk
embeddings. Jika informasi belum ada di dokumen, chatbot dapat memberikan
jawaban fallback atau meminta eskalasi ke HR.}

\Needspace{12\baselineskip}
\subsection{Riwayat Chat}

Percakapan tersimpan per akun dan muncul pada daftar \textbf{Chats} di sidebar.

\begin{itemize}
  \item Klik salah satu judul untuk membuka kembali percakapan lama.
  \item Klik \textbf{New chat} di kanan atas untuk memulai percakapan baru.
  \item Setiap baris punya aksi untuk mengganti nama dan menghapus percakapan.
\end{itemize}

\notebox{Riwayat bersifat pribadi per akun: Anda hanya melihat percakapan Anda
sendiri. Untuk melihat percakapan seluruh karyawan, gunakan menu \textbf{Logs}.}

\Needspace{16\baselineskip}
\section{Feedback Jawaban Chat}

Setiap jawaban asisten menampilkan baris \textbf{Is this helpful?} dengan dua
ikon jempol tepat di bawah info \textit{AI Assistant}, waktu, dan sumber jawaban.

\screenshot[0.66\linewidth]{17-chat-feedback-row.png}{Baris feedback di bawah jawaban asisten.}

\begin{center}
\renewcommand{\arraystretch}{1.2}
\begin{tabularx}{0.92\linewidth}{@{}B{3.2cm}Y@{}}
\toprule
Ikon & Fungsi \\
\midrule
Jempol atas (hijau) & Menandai jawaban itu bagus. Klik langsung tersimpan tanpa perlu mengisi alasan. \\
Jempol bawah (merah) & Melaporkan jawaban yang kurang memuaskan. Setelah diklik, muncul form kecil untuk mengisi alasan sebelum feedback terkirim. \\
\bottomrule
\end{tabularx}
\end{center}

\notebox{Jumlah feedback negatif ikut terhitung pada kartu \textbf{Feedback} di
menu Logs dan KPI \textbf{Negative Feedback} di menu Analytics, jadi feedback
dari user memang dipakai untuk memantau kualitas jawaban.}

\Needspace{18\baselineskip}
\section{FAQ}

Fitur FAQ dipakai untuk membuat daftar pertanyaan umum berdasarkan dokumen yang
sudah masuk knowledge base.

\Needspace{14\baselineskip}
\subsection{Generate FAQ}

\begin{enumerate}
  \item Buka menu \textbf{FAQ}.
  \item Isi pertanyaan pada kolom \textbf{Question}.
  \item Klik \textbf{Generate FAQ}.
\end{enumerate}

\screenshot[0.76\linewidth]{05-faq-generate.png}{Panel admin untuk generate FAQ.}

\Needspace{14\baselineskip}
\subsection{Update dan Delete FAQ}

FAQ yang sudah dibuat dapat diperbarui atau dihapus dari daftar.

\begin{itemize}
  \item Klik \textbf{Update} untuk mengubah pertanyaan dan generate ulang jawabannya.
  \item Klik \textbf{Delete} untuk menghapus FAQ.
\end{itemize}

\screenshot[0.55\linewidth]{07-faq-update-delete.png}{Tombol Update dan Delete pada item FAQ.}

Saat update, pertanyaan akan masuk kembali ke form dan tombol berubah menjadi
\textbf{Regenerate FAQ}.

\screenshot[0.52\linewidth]{08-faq-regenerate.png}{Mode regenerate FAQ.}

\Needspace{18\baselineskip}
\section{Document Library}

Menu \textbf{Docs} dipakai untuk mengelola dokumen SOP dan Form yang tersedia di
aplikasi.

\screenshot[0.76\linewidth]{09-document-management.png}{Panel admin untuk mengelola dokumen.}

\Needspace{12\baselineskip}
\subsection{Mencari Dokumen}

Klik tombol \textbf{Filter} di kanan atas untuk memunculkan kolom pencarian,
lalu ketik sebagian nama dokumen untuk menyaring daftar.

\Needspace{12\baselineskip}
\subsection{Upload / Insert Dokumen}

\begin{enumerate}
  \item Klik \textbf{Choose files}.
  \item Pilih satu atau beberapa file \texttt{pdf}, \texttt{docx}, atau \texttt{txt}.
  \item Klik \textbf{Insert files}.
\end{enumerate}

\Needspace{12\baselineskip}
\subsection{Update dan Delete Dokumen}

Setiap baris dokumen memiliki tombol \textbf{Download}, \textbf{Update}, dan
\textbf{Delete}.

\screenshot[0.52\linewidth]{10-document-row-actions.png}{Aksi pada baris dokumen.}

\Needspace{14\baselineskip}
\subsection{Menambahkan Form ke SOP}

Form (Word/Excel/PDF) ditautkan langsung ke satu SOP.

\begin{enumerate}
  \item Klik tombol \textbf{Forms} pada baris SOP yang dituju untuk membuka panel form terkait.
  \item Klik \textbf{Insert form} (atau \textbf{Insert another form} jika sudah ada form sebelumnya).
  \item Pilih file form, lalu tunggu proses upload selesai.
\end{enumerate}

\screenshot[0.62\linewidth]{18-document-insert-form.png}{Panel Forms pada baris SOP dengan tombol Insert another form.}

\notebox{Saat form diunduh, aplikasi menanyakan format terlebih dahulu lewat
dialog \textbf{Choose format} dengan pilihan \textbf{PDF} atau \textbf{Word}.}

\Needspace{12\baselineskip}
\subsection{Finalize Perubahan}

Klik \textbf{Finalize} setelah perubahan pada SOP/knowledge document agar chatbot
memakai sumber terbaru.

\begin{enumerate}
  \item Insert, update, atau delete dokumen SOP.
  \item Klik \textbf{Finalize}.
  \item Tunggu proses selesai (Akan lebih lama jika ada banyak perubahan, bisa pindah ke pekerjaan lain sambil menunggu).
  \item Tes pertanyaan terkait SOP tersebut melalui chat.
\end{enumerate}

\screenshot[0.52\linewidth]{11-rebuild-embeddings.png}{Tombol Finalize aktif setelah SOP diperbarui.}

\notebox{Selama perubahan belum di-Finalize, tombol \textbf{Undo all} tersedia
untuk membatalkan seluruh perubahan dokumen yang belum diterapkan.}

\Needspace{20\baselineskip}
\section{Logs}

Menu \textbf{Logs} digunakan untuk memantau penggunaan chatbot oleh seluruh
karyawan. Di bagian atas ada tiga kartu ringkasan yang sekaligus berfungsi
sebagai pemilih tampilan: klik salah satu kartu untuk mengganti isi daftar di
bawahnya.

\begin{center}
\renewcommand{\arraystretch}{1.2}
\begin{tabularx}{0.94\linewidth}{@{}B{3.4cm}Y@{}}
\toprule
Kartu & Isi daftar setelah diklik \\
\midrule
Total Chat & \textbf{Recent questions} --- seluruh pertanyaan pada rentang tanggal terpilih. \\
Total Session & \textbf{Chat sessions} --- percakapan dikelompokkan per sesi, lengkap dengan jumlah pertanyaan dan nama user. \\
Feedback & Daftar jawaban yang dilaporkan user sebagai tidak memuaskan. \\
\bottomrule
\end{tabularx}
\end{center}

\screenshot[0.72\linewidth]{12-logs-questions-detail.png}{Daftar pertanyaan; satu baris dibuka untuk melihat jawaban chatbot.}

Pada toolbar tersedia beberapa alat bantu:

\begin{itemize}
  \item \textbf{Search by name} --- menyaring aktivitas berdasarkan nama karyawan.
  \item \textbf{Start} dan \textbf{End} --- rentang tanggal yang ditampilkan.
  \item Tombol \textbf{refresh} --- memuat ulang data terbaru.
\end{itemize}

Klik satu baris pertanyaan untuk membuka jawaban yang diberikan chatbot. Jika
jawabannya panjang, gunakan tombol \textbf{Read more} untuk melihat selengkapnya.

\screenshot[0.72\linewidth]{13-logs-sessions.png}{Tampilan Chat sessions setelah kartu Total Session diklik.}

Saat satu sesi dibuka, filter \textbf{Filtered session} akan aktif dan daftar
hanya menampilkan pertanyaan dari sesi tersebut. Klik \textbf{Clear} untuk
kembali melihat semua pertanyaan.

\screenshot[0.72\linewidth]{14-logs-session-filter.png}{Log sesi terpilih dengan tombol Clear dan detail jawaban.}

\notebox{Tombol tempat sampah pada setiap baris menghapus log tersebut secara
permanen. Data yang sudah dihapus tidak bisa dikembalikan dan akan hilang juga
dari perhitungan Analytics.}

\Needspace{20\baselineskip}
\section{Analytics}

Menu \textbf{Analytics} menampilkan ringkasan pemakaian chatbot: seberapa ramai
dipakai, topik apa yang paling banyak ditanyakan, dan seberapa sering jawaban
dinilai kurang memuaskan.

\screenshot[0.78\linewidth]{19-analytics-overview.png}{Halaman Usage Analytics dengan pemilih rentang tanggal dan tiga KPI utama.}

\Needspace{12\baselineskip}
\subsection{Memilih Rentang Tanggal}

Pada bar \textbf{Overview} tersedia tombol cepat \textbf{7D}, \textbf{30D}, dan
\textbf{90D}, atau isi sendiri kolom \textbf{Start} dan \textbf{End}. Seluruh
angka dan grafik di halaman ini mengikuti rentang yang dipilih.

Tombol \textbf{refresh} di sebelah kanan menghitung ulang data agregat di server.
Keterangan \textbf{Last refreshed} dan \textbf{Data available} di kanan atas
menunjukkan kapan angka terakhir dihitung dan periode data yang tersedia.

\Needspace{14\baselineskip}
\subsection{Tiga KPI Utama}

\begin{center}
\renewcommand{\arraystretch}{1.2}
\begin{tabularx}{0.94\linewidth}{@{}B{3.9cm}Y@{}}
\toprule
KPI & Arti \\
\midrule
Total Questions & Jumlah pertanyaan pada rentang terpilih. \\
Active Users & Jumlah karyawan berbeda yang bertanya pada rentang terpilih. \\
Negative Feedback & Jumlah jawaban yang dilaporkan tidak memuaskan, beserta persentasenya terhadap total pertanyaan. \\
\bottomrule
\end{tabularx}
\end{center}

Di bawah setiap KPI ada perbandingan terhadap periode sebelumnya dengan panjang
yang sama. Jika belum ada data pembanding, tertulis \textit{No prior-period data}.

\Needspace{16\baselineskip}
\subsection{Grafik dan Tabel}

\screenshot[0.78\linewidth]{20-analytics-charts.png}{Grafik tren, sebaran topik, tabel topik, dan daftar user paling aktif.}

\begin{center}
\renewcommand{\arraystretch}{1.2}
\begin{tabularx}{0.94\linewidth}{@{}B{3.9cm}Y@{}}
\toprule
Panel & Isi \\
\midrule
Questions over time & Jumlah pertanyaan per hari, untuk melihat hari tersibuk. \\
Topic share & Proporsi pertanyaan per topik SOP dalam bentuk donat beserta persentasenya. \\
Topic performance & Banyaknya karyawan berbeda yang menanyakan tiap topik, diurutkan dari yang paling banyak. \\
Most active users & Karyawan dengan jumlah pertanyaan terbanyak pada rentang terpilih. \\
\bottomrule
\end{tabularx}
\end{center}

\notebox{Panel \textbf{Topic share}, \textbf{Topic performance}, dan \textbf{Most
active users} bisa diklik. Klik satu bagian donat, satu baris topik, atau satu
baris user untuk langsung membuka menu \textbf{Logs} yang sudah tersaring sesuai
pilihan tersebut.}

\Needspace{20\baselineskip}
\section{Guardrails}

Menu \textbf{Guardrails} dipakai untuk mengatur batasan yang disuntikkan ke
setiap jawaban AI: topik apa saja yang boleh dijawab, dan bagaimana chatbot harus
bersikap saat pertanyaan di luar topik SOP atau saat ada percobaan
mengubah/membocorkan instruksi sistemnya.

\screenshot[0.78\linewidth]{15-guardrails-menu.png}{Menu Guardrails dengan daftar rules yang sedang aktif.}

\begin{enumerate}
  \item Ubah teks aturan pada kolom editor.
  \item Klik \textbf{Save guardrails}.
  \item Tes lewat chat untuk memastikan chatbot sudah mengikuti aturan baru.
\end{enumerate}

\notebox{Aturan di sini ikut terkirim pada setiap jawaban, jadi perubahan
langsung terasa di semua percakapan. Ubah seperlunya dan pastikan hasilnya
diuji.}

\Needspace{14\baselineskip}
\section{Logout}

Klik ikon logout pada panel akun di kiri bawah, lalu konfirmasi dengan tombol
\textbf{Logout}.

\screenshot[0.44\linewidth]{23-logout.png}{Konfirmasi logout.}

\notebox{Setelah logout, riwayat chat pada layar dibersihkan dan aplikasi kembali
tertutup halaman sign-in. Selalu logout jika memakai komputer bersama.}

\Needspace{12\baselineskip}
\section{Checklist Admin}

\begin{center}
\setlength{\fboxsep}{0.7em}
\fcolorbox{icsgray}{icsnote}{
  \begin{minipage}{0.84\linewidth}
    \textbf{Checklist Operasional Admin}
    \vspace{0.4em}

    \begin{tabularx}{\linewidth}{@{}p{0.8cm}Y@{}}
      \checkline{Login memakai akun Google Workspace \textbf{@icscompute.com}.}
      \checkline{Tambahkan admin baru cukup dengan email kantornya, tanpa password.}
      \checkline{Setelah mengubah SOP, klik \textbf{Finalize}.}
      \checkline{Tambahkan Form lewat panel \textbf{Forms} pada baris SOP terkait, bukan lewat \textbf{Insert files}.}
      \checkline{Cek \textbf{Logs} untuk melihat pertanyaan user, sesi chat, dan feedback negatif.}
      \checkline{Cek \textbf{Analytics} untuk melihat tren pemakaian dan topik yang paling banyak ditanyakan.}
      \checkline{Logout setelah selesai, terutama di komputer bersama.}
    \end{tabularx}
  \end{minipage}
}
\end{center}

\end{document}
